\section{Kummer duality and the group input}\label{sec:preliminaries}

We first give the topological form of Kummer duality needed for selected
subfields and residue fields.  In this section the residue characteristic
is allowed to equal $\ell$ unless explicitly excluded.

For a function field $F/\kappa$ over algebraically closed constants, put
\[
 W_F=\Hom(F^\times,\Zell),\qquad
 \widehat F=\varprojlim_{e\ge1} F^\times/F^{\times\ell^e}.
\]
Equip $W_F$ with pointwise convergence and $\widehat F$ with its
$\ell$-adic topology.  We use additive notation for these modules and
write $[f]$ for the class of $f\in F^\times$.  In particular
$[fg]=[f]+[g]$.  All characters vanish on $\kappa^\times$.  We abbreviate
$F^\times/F^{\times q}$ to $F^\times/q$.  An annihilator $B^\perp$
of $B\subset F^\times$ or $B\subset\widehat F$ is taken in $W_F$;
for $B\subset W_F$ it is taken in $\widehat F$, using evaluation.

\begin{lemma}[Continuous evaluation]\label{prelim:duality}
The group $F^\times/\kappa^\times$ is free abelian.  Consequently
$W_F$ is a product of copies of $\Zell$, and evaluation gives
\[
 \widehat F=\Hom_{\mathrm{cont}}(W_F,\Zell),\qquad
 \widehat F/\ell^e\widehat F=F^\times/F^{\times\ell^e}.
\]
The module $\widehat F$ is torsion free, and evaluations on $W_F$
detect every finite-level class.
\end{lemma}
\begin{proof}
Choose a normal integral projective model of $F/\kappa$.  A rational
function with zero Weil divisor on this model is a global unit and hence
a constant.  Thus the divisor map embeds $F^\times/\kappa^\times$ in
the free abelian group of prime divisors; a subgroup of a free abelian
group is free.

For a basis indexed by a set $J$, the character group is $\Zell^J$.
Modulo $\ell^e$, a continuous character on this compact product depends
on finitely many coordinates.  Compatible such characters are therefore
families $(c_j)_{j\in J}$ in $\Zell$ with only finitely many
$c_j\notin\ell^e\Zell$ for each $e$.  This is exactly the completion
of $\bigoplus_{j\in J}\mathbb Z$.  Reduction modulo $\ell^e$ and
coordinate evaluations prove the assertions.  Divisibility of
$\kappa^\times$ removes the constants at every level.
\end{proof}

When $\operatorname{char}F\ne\ell$, a choice of a compatible system of
$\ell$-power roots of unity identifies the Tate module $\Zell(1)$ with
$\Zell$.  Kummer theory then identifies $\Pi_F^a$ with $W_F$.
An isomorphism $\varphi:\Pi_L^a\to\Pi_K^a$ has evaluation dual
\begin{equation}\label{eq:theta}
 \Theta:\widehat K\longrightarrow\widehat L,
 \qquad \Theta(h)(\sigma)=h(\varphi\sigma).
\end{equation}
Changing the Tate identifications multiplies this map by one unit.
We fix the identifications for the proof.

\begin{lemma}[Subfields and finite extensions]\label{prelim:completion}
Let $E\subset F$ be function fields over the same algebraically closed
field $\kappa$.
\begin{enumerate}
\item If $E$ is relatively algebraically closed in $F$, then
$\widehat E\to\widehat F$ is injective, its image is closed and
$\ell$-saturated, and restriction $W_F\to W_E$ is a continuous
surjection.  Here $\ell$-saturated means that $\ell h\in\widehat E$
for $h\in\widehat F$ implies $h\in\widehat E$.
\item If $F/E$ is finite, its completion map is injective, and
inclusion followed by norm on completions is multiplication by
$[F:E]$.
\end{enumerate}
\end{lemma}
\begin{proof}
In the first case, a root in $F$ of an element of $E$ is algebraic over
$E$, so injectivity holds at every finite level.  Alternatively, on a
normal integral projective $E$-model of $F$, the divisor map embeds
$F^\times/E^\times$ into a free abelian group: its kernel consists of
global units, which are algebraic over $E$ and therefore lie in $E$.
Thus $F^\times/E^\times$ is free abelian, and
\[
 0\longrightarrow E^\times/\kappa^\times
 \longrightarrow F^\times/\kappa^\times
 \longrightarrow F^\times/E^\times\longrightarrow0
\]
splits as a sequence of abelian groups.  Completion preserves this
splitting.  The two completed free factors are torsion free, proving
closedness and saturation.  Every character on $E^\times$ extends along
the splitting, proving the asserted surjection; continuity follows from
the pointwise topologies.

For a finite extension, the field norm is defined also in the inseparable
case and sends an element of $E$ to its $[F:E]$-th power.  It induces
the stated composite on completions.  Torsion freeness from
Lemma~\ref{prelim:duality} proves injectivity.
\end{proof}

The next observation lets us apply Pop's reconstruction theorems to
the stated group input.

\begin{lemma}[Lifting bracket-compatible isomorphisms]
\label{prelim:central-lift}
Every bracket-compatible isomorphism in Theorem~\ref{thm:main} is the
abelianization of a continuous isomorphism $\Pi_L^c\to\Pi_K^c$.
This holds also for $\ell=2$.
\end{lemma}
\begin{proof}
Use the given maps to identify the bases of the two central extensions
with one abelian pro-$\ell$ group $A$ and their kernels with one group
$B$.  Let $E_1,E_2$ denote these extensions.  In the fiber product
$P=E_1\times_A E_2$, quotient by the diagonal central subgroup
$\{(b,b):b\in B\}$.  The resulting group $Q$ is abelian: the
commutator of two pairs has identical components by bracket
compatibility.  It fits into an exact sequence of abelian pro-$\ell$
groups
\[
 0\longrightarrow B\longrightarrow Q\longrightarrow A\longrightarrow0.
\]
By Lemma~\ref{prelim:duality}, $A$ is a product of copies of $\Zell$.
Such products are projective in the category of abelian pro-$\ell$
groups.  Indeed Pontryagin duality changes $A$ into a direct sum of
copies of $\mathbb Q_\ell/\Zell$, a divisible, hence injective,
discrete torsion group.  Consequently $Q\to A$ has a continuous
section.

The inverse image of this section in $P$ maps isomorphically to each
$E_i$.  For fixed first component, the second component can be adjusted
by a unique element of $B$ to land over the section; its uniqueness
also proves that each projection has trivial kernel.  These are
continuous bijections between compact Hausdorff groups, hence
isomorphisms.  The resulting map $E_1\to E_2$ induces the prescribed
base and kernel maps.  The proof uses no division by $2$ and leaves
no additional power datum to specify.
\end{proof}

\subsection{Orders on a curve}
For a smooth projective integral curve $C/\kappa$ with function field
$E$, its closed points are $\kappa$-rational.  Each normalized order
$\ord_P:E^\times\to\mathbb Z$ extends to $\widehat E\to\Zell$.

\begin{lemma}[Curve kernels]\label{prelim:curve-kernel}
If $\operatorname{char}\kappa\ne\ell$, then
\[
 \ker\left(E^\times/E^{\times\ell}
       \xrightarrow{(\ord_P)_P}\bigoplus_{P\in C}\mathbb Z/\ell\right)
 \simeq\Pic(C)[\ell].
\]
Its dimension over $\Fell$ is $2g(C)$.  The kernel of all exact point
orders on $\widehat E$ is the Tate module $T_\ell\Pic(C)$ and has
$\Zell$-rank $2g(C)$.
\end{lemma}
\begin{proof}
The same argument works with $q=\ell^e$.  If
$\operatorname{div}(f)=qD$, send $[f]$ to the divisor class of $D$.
Changing $f$ by a $q$-th power does not change this class.  Every
$q$-torsion divisor class occurs.  The kernel consists of functions
$cu^q$; the constant $c\in\kappa^\times$ is a $q$-th power.  This
proves the finite-level identification.  The $q$-torsion of the
Jacobian over an algebraically closed field of characteristic prime to
$q$ is $(\mathbb Z/q)^{2g(C)}$
\cite[Tag~0C1Z, Lemma~53.17.1]{Stacks}.  Under reduction from $\ell^{e+1}$
to $\ell^e$, the divisor class is multiplied by $\ell$.  Taking the
inverse limit gives the exact-order kernel and its rank.
\end{proof}

In particular, an isomorphism of the $W$-groups carrying the family of
point inertia lines bijectively onto the corresponding family preserves
the genus.  A generator
of each line can change by a unit; this changes neither vanishing nor
nonvanishing of any finite-level order.  We will use this observation
only after proving that the residue characteristics differ from $\ell$.
